\documentclass[11pt, letterpaper]{article}
\usepackage{mobi-lectura}

\title{Sistemas de Colas con SED:\\[4pt]
Ley de Little, Período Transitorio y Validación}
\author{Alejandra Tabares Pozos}
\date{}

\begin{document}

\maketitle
\tableofcontents
\newpage

% ══════════════════════════════════════════════════
\section{Introducción}
% ══════════════════════════════════════════════════

La lectura anterior (Mecánica, Trazas y Worldviews) mostró cómo funciona internamente un motor de SED: el avance al siguiente evento, la FEL y las rutinas de llegada y salida. Las trazas realizadas con 6--7 clientes ilustraron la mecánica, pero no reflejan el comportamiento real de un sistema en operación continua.

Esta lectura aborda las tres preguntas que surgen al pasar de la traza manual al estudio de simulación:
\begin{enumerate}
    \item ¿Qué relación hay entre el número de entidades en el sistema y el tiempo que pasan en él? (\textbf{Ley de Little})
    \item ¿Cuánto tiempo debe correr la simulación antes de que las estadísticas sean confiables? (\textbf{Período transitorio})
    \item ¿Cómo sé que mis resultados son correctos? (\textbf{Validación y buenas prácticas})
\end{enumerate}

% ══════════════════════════════════════════════════
\section{Simulación de colas en estado estable: el sistema $M/M/2$}
% ══════════════════════════════════════════════════

Consideremos un sistema $M/M/2$ con tasa de llegadas $\lambda = 5$ clientes/minuto, tasa de servicio $\mu = 3$ por servidor y $c = 2$ servidores idénticos. La intensidad de tráfico es $\rho = \lambda/(c\mu) = 5/6 \approx 0{,}833 < 1$, lo que garantiza estabilidad.

\subsection{Métricas analíticas de referencia}

Para $M/M/c$ existen fórmulas cerradas. La probabilidad de estado vacío es $P_0 \approx 0{,}0909$ y la fórmula C de Erlang da la probabilidad de esperar: $C(2, 5/3) \approx 0{,}7576$. Las métricas de estado estable son:

\begin{center}
\renewcommand{\arraystretch}{1.4}
\small
\begin{tabular}{@{}lcc@{}}
\toprule
\textbf{Métrica} & \textbf{Fórmula} & \textbf{Valor} \\
\midrule
$L_q$ (clientes en cola) & $C(c,\lambda/\mu)\cdot\rho/(1-\rho)$ & 3{,}788 \\
$L$ (clientes en sistema) & $L_q + \lambda/\mu$ & 5{,}455 \\
$W_q$ (espera en cola) & $L_q/\lambda$ & 0{,}758 min \\
$W$ (tiempo en sistema) & $W_q + 1/\mu$ & 1{,}091 min \\
$\rho$ (utilización) & $\lambda/(c\mu)$ & 0{,}833 \\
\bottomrule
\end{tabular}
\end{center}

Estos valores corresponden al \textbf{estado estable} ($t \to \infty$). La simulación debe converger hacia ellos una vez superado el período transitorio.

\subsection{Ejemplo paso a paso con 10 clientes}

Al simular los primeros 10 clientes con tiempos generados a partir de las distribuciones exponenciales correspondientes, se observa que las esperas crecen progresivamente: los clientes 1 y 2 no esperan (llegan a un sistema vacío), pero a partir del cliente 3 las demoras aumentan de 0{,}03 a 0{,}44 minutos. El sistema está en fase de \emph{llenado} (transitorio). Las métricas parciales en $T = 2$ minutos arrojan $\hat{W}_D = 0{,}579$ y $\hat{W}_{q,D} = 0{,}182$, con errores del 47\% y 76\% respecto a los valores teóricos. Esta discrepancia ilustra de manera contundente por qué no podemos confiar en estimaciones basadas en corridas cortas.


% ══════════════════════════════════════════════════
\section{La Ley de Little}
% ══════════════════════════════════════════════════

\begin{teorema}[Ley de Little, 1961]
En estado estable, para cualquier sistema de colas:
\[
\boxed{L = \lambda \cdot W}
\]
donde $L$ es el número promedio de entidades en el sistema, $\lambda$ es la tasa de llegada efectiva y $W$ es el tiempo promedio en el sistema.
\end{teorema}

La ley se aplica a cualquier subsistema: $L_q = \lambda \cdot W_q$ (para la cola) y $L_s = \lambda \cdot W_s$ (para el servicio). Su generalidad es extraordinaria: no depende de la distribución de llegadas, de la distribución de servicio, del número de servidores, de la disciplina de cola ni del orden de servicio. La única condición fundamental es que el sistema esté en estado estable.

\subsection{Tres versiones de Little}

Es crucial distinguir tres versiones de la ley, pues su ámbito de validez difiere:

La \textbf{versión estocástica} $L = \lambda W$ es un resultado límite que requiere estado estable. La \textbf{versión operacional completa}, $\bar{L}(T) = \bar{\lambda}(T) \cdot \bar{W}_A(T)$, donde $\bar{W}_A$ incluye a \emph{todos} los clientes (incluso los que siguen en el sistema al tiempo $T$, con su tiempo parcial), es una \textbf{identidad contable} que se cumple siempre, para todo $T$, incluso en el transitorio. La demostración es inmediata: el área bajo $N(t)$ equivale a la suma de los tiempos individuales en sistema. La \textbf{versión operacional parcial}, $\bar{L}(T) \approx \bar{\lambda}(T) \cdot \bar{W}_D(T)$, que solo promedia sobre clientes que ya completaron servicio, \emph{no} coincide con $\bar{L}$ durante el transitorio porque $\bar{W}_D$ ignora a los clientes aún presentes en el sistema.

En la práctica, esto significa que si usamos $\bar{W}_D$ como estimador (lo más común), debemos descartar el período transitorio para que la relación de Little sea consistente.

\subsection{Little como herramienta de verificación}

La Ley de Little es una de las herramientas más poderosas para \textbf{validar} un simulador. Si al finalizar una corrida larga las relaciones $\hat{L} \approx \hat{\lambda}\cdot\hat{W}$ y $\hat{L}_q \approx \hat{\lambda}\cdot\hat{W}_q$ no se cumplen razonablemente, algo está mal: el período transitorio no fue eliminado, la simulación fue demasiado corta, hay un bug en el código (contadores mal actualizados), el $\bar{W}$ se calculó solo con completados sin descartar el warm-up, o el sistema es inestable ($\rho \geq 1$).


% ══════════════════════════════════════════════════
\section{El período transitorio y el estado estable}
% ══════════════════════════════════════════════════

\subsection{Convergencia al estado estable}

Cuando un sistema estable ($\rho < 1$) arranca desde una condición inicial (típicamente vacío), sus métricas no alcanzan inmediatamente los valores de estado estable. Durante el \textbf{período transitorio} la cola está en fase de llenado, las esperas crecen y las estimaciones están sesgadas.

La siguiente figura ilustra cómo $\bar{L}(T)$ converge al valor teórico $L^* = 5{,}455$ para el $M/M/2$. Al inicio, la estimación es mucho menor que $L^*$ porque el sistema arranca vacío.

\begin{figure}[H]
\centering
\begin{tikzpicture}
\begin{axis}[
    width=13cm, height=5.5cm,
    xlabel={Tiempo de simulación $T$},
    ylabel={$\bar{L}(T)$},
    xmin=0, xmax=100,
    ymin=0, ymax=9,
    grid=both, grid style={gray!20},
    legend pos=north east,
    legend style={font=\small},
]
\addplot[thick, uniandes, domain=0:100, samples=300]
    {5.4545 + 4*exp(-0.05*x)*cos(deg(0.15*x)) - 5.4545*exp(-0.08*x)};
\addlegendentry{$\bar{L}(T)$ observado}

\addplot[dashed, black, domain=0:100] {5.4545};
\addlegendentry{$L^* = 5.455$}

\addplot[dashed, gray!60, domain=0:100] {5.4545 + 0.3};
\addplot[dashed, gray!60, domain=0:100] {5.4545 - 0.3};

\draw[<->, thick, rojoclaro] (axis cs:0,8.5) -- (axis cs:40,8.5);
\node[above, font=\small] at (axis cs:20,8.5) {\textcolor{rojoclaro}{\textbf{Transitorio}}};

\draw[<->, thick, verdeoscuro] (axis cs:40,8.5) -- (axis cs:100,8.5);
\node[above, font=\small] at (axis cs:70,8.5) {\textcolor{verdeoscuro}{\textbf{Estado estable}}};
\end{axis}
\end{tikzpicture}
\caption{Convergencia de $\bar{L}(T)$ al valor teórico para el $M/M/2$ con $\rho = 0.833$.}
\label{fig:convergencia}
\end{figure}

\subsection{Efecto de las condiciones iniciales}

Si el sistema arranca vacío, $\bar{L}(T)$ subestima $L^*$; si arranca sobrecargado, la sobreestima. Ambas curvas convergen al mismo valor. Esto muestra que el sesgo es una consecuencia del punto de partida, no del simulador.

\subsection{Sensibilidad a $\rho$}

A mayor $\rho$, más largo es el transitorio. La siguiente figura compara tres niveles de carga para el mismo sistema.

\begin{figure}[H]
\centering
\begin{tikzpicture}
\begin{axis}[
    width=13cm, height=5.5cm,
    xlabel={Tiempo de simulación},
    ylabel={$\bar{L}(T) / L^*$},
    xmin=0, xmax=150,
    ymin=0, ymax=1.3,
    grid=both, grid style={gray!20},
    legend pos=south east,
    legend style={font=\small},
]
\addplot[thick, verdeoscuro, domain=0:150, samples=200]
    {1 - exp(-0.15*x)};
\addlegendentry{$\rho = 0.50$}

\addplot[thick, uniandes, domain=0:150, samples=200]
    {1 - exp(-0.06*x)};
\addlegendentry{$\rho = 0.83$}

\addplot[thick, rojoclaro, domain=0:150, samples=200]
    {1 - exp(-0.02*x)};
\addlegendentry{$\rho = 0.95$}

\addplot[dashed, black, domain=0:150] {1};

\draw[dashed, verdeoscuro] (axis cs:20,0) -- (axis cs:20,1.3);
\draw[dashed, uniandes] (axis cs:50,0) -- (axis cs:50,1.3);
\draw[dashed, rojoclaro] (axis cs:120,0) -- (axis cs:120,1.3);
\end{axis}
\end{tikzpicture}
\caption{Sistemas muy cargados ($\rho = 0.95$) requieren transitorios mucho más largos.}
\label{fig:sensibilidad}
\end{figure}

\subsection{Método de Welch}

El procedimiento de Welch es el más utilizado para detectar el transitorio. Consiste en: (1)~ejecutar $R$ réplicas independientes (típicamente $R \geq 10$); (2)~para cada réplica $r$, registrar la trayectoria de interés $Y_r(t)$; (3)~calcular el promedio entre réplicas $\bar{Y}(t) = \frac{1}{R}\sum_r Y_r(t)$; (4)~aplicar una media móvil $\hat{Y}(t) = \frac{1}{2w+1}\sum_{s=t-w}^{t+w}\bar{Y}(s)$ para suavizar; (5)~identificar visualmente el tiempo $T_0$ a partir del cual $\hat{Y}(t)$ se estabiliza.

\subsection{Estrategias para manejar el transitorio}

Existen varias estrategias. El \textbf{truncamiento} (\emph{warm-up deletion}) es la más común: se descarta toda la información de $[0, T_0]$ y se recolectan estadísticas solo en $[T_0, T]$. Una \textbf{corrida larga} consiste en hacer $T$ tan grande que el transitorio sea despreciable respecto al total. El \textbf{inicio inteligente} arranca el sistema con un estado cercano al estable (si se conoce $L^*$ aproximadamente). Las \textbf{réplicas con truncamiento} combinan $R$ réplicas independientes, cada una con su warm-up, para construir intervalos de confianza. Finalmente, el \textbf{método de batch means} divide una única corrida larga en lotes después de $T_0$.

Para el sistema $M/M/2$ con $\rho = 0{,}833$, el transitorio dura entre 20 y 50 unidades de tiempo. Sin truncamiento, la estimación de $L$ tiene un sesgo del $-12\%$; con $T_0 = 30$ minutos, el sesgo se reduce a menos del $1\%$.


% ══════════════════════════════════════════════════
\section{Buenas prácticas en simulación de colas}
% ══════════════════════════════════════════════════

A modo de resumen operativo, todo estudio de simulación robusto debe seguir un protocolo sistemático:

\begin{enumerate}
    \item \textbf{Verificar la estabilidad}: asegurar $\rho < 1$ antes de buscar métricas de estado estable.
    \item \textbf{Detectar el transitorio}: usar el método de Welch u otro procedimiento formal; no asumir un $T_0$ arbitrario.
    \item \textbf{Verificar Little}: comprobar que $\hat{L} \approx \hat{\lambda}\cdot\hat{W}$ como diagnóstico de consistencia.
    \item \textbf{Usar suficientes réplicas}: un mínimo de $R \geq 10$ para construir intervalos de confianza.
    \item \textbf{Corridas suficientemente largas}: $T \gg T_0$ (al menos $T/T_0 > 10$).
    \item \textbf{Validar contra lo conocido}: si existe solución analítica, comparar.
    \item \textbf{Reportar incertidumbre}: acompañar estimaciones puntuales con intervalos de confianza.
\end{enumerate}


% ══════════════════════════════════════════════════
\section{Conclusiones}
% ══════════════════════════════════════════════════

La Ley de Little, $L = \lambda W$, es el resultado más fundamental de la teoría de colas. En su versión operacional completa es una identidad contable que siempre se cumple; en su versión estocástica requiere estado estable, lo que obliga a detectar y eliminar el período transitorio. Ignorar el transitorio puede producir errores de estimación superiores al 10\%, y a mayor $\rho$ el problema se agrava.

Las métricas del sistema se calculan de dos formas complementarias: promedios en el tiempo (áreas bajo trayectorias) y promedios por entidad (sumas sobre clientes). Dominar ambas perspectivas y saber cuándo cada una es apropiada es esencial para el análisis correcto de la salida de una simulación.

\vfill

\section*{Referencias}

\begin{enumerate}[label={[\arabic*]}]
    \item Banks, J., Carson, J.\,S., Nelson, B.\,L., Nicol, D.\,M. (2014). \textit{Discrete-Event System Simulation}, 5th ed. Pearson. Cap.~9.
    \item Law, A.\,M. (2015). \textit{Simulation Modeling and Analysis}, 5th ed. McGraw-Hill. Cap.~9.
    \item Little, J.\,D.\,C. (1961). A proof for the queuing formula: $L = \lambda W$. \textit{Operations Research}, 9(3), 383--387.
    \item Stidham, S. (1974). A last word on $L = \lambda W$. \textit{Operations Research}, 22(2), 417--421.
    \item Welch, P.\,D. (1983). The statistical analysis of simulation results. In \textit{The Computer Performance Modeling Handbook}, Academic Press.
    \item Gross, D., Shortle, J., Thompson, J., \& Harris, C. (2008). \textit{Fundamentals of Queueing Theory}, 4th ed. Wiley. Cap.~2.
\end{enumerate}

\end{document}
